\section{Related Work}
\label{sec:related_work}

Our work is situated at the intersection of Parameter-Efficient Fine-Tuning, Model Merging, and Post-Training Quantization. In this section, we review relevant literature in these areas and highlight the methodological gaps we address.

\subsection{Parameter-Efficient Fine-Tuning (PEFT)}
As pre-trained base models have scaled, fine-tuning all model parameters has become prohibitively expensive. This has spurred the development of Parameter-Efficient Fine-Tuning (PEFT) methods, with Low-Rank Adaptation (LoRA) \cite{Hu2021LoRA} being the most prominent. LoRA factorizes the weight updates into low-rank matrices $A$ and $B$, freezing the pre-trained weights and training only these small matrices. QLoRA \cite{Dettmers2023QLoRA} further extends this by quantizing the pre-trained base model to NormalFloat4 (NF4) and introducing double quantization, backpropagating gradients through dequantized weights to update the high-precision LoRA parameters. While QLoRA achieves near-lossless performance, its standard pipeline assumes that the high-precision adapters can be seamlessly merged back into the base model and deployed. We critically evaluate this assumption.

\subsection{Model Merging}
Model merging seeks to combine multiple, independently fine-tuned models of identical architecture into a single model without additional training. Early work focused on simple weight averaging \cite{Wortsman2022ModelSoups} or permutation alignment to resolve coordinate-matching issues \cite{Ainsworth2022GitReBasin,Stoica2023ZipIt}. More recently, task-vector-based approaches have achieved prominence. Task Arithmetic \cite{Ilharco2022TaskArithmetic} creates task vectors by subtracting pre-trained weights from fine-tuned weights and scaling them before addition. TIES-Merging \cite{Yadav2023TIES} resolves parameter interference by keeping only the top-valued components and resolving sign conflicts. DARE \cite{Yu2023DARE} randomly drops small parameter values to sparsify updates. Other approaches optimize layer-wise merging coefficients via test-time adaptation \cite{Wang2021Tent} or entropy minimization, such as AdaMerging \cite{MergingSub3} and RegCalMerge \cite{MergingSub3}. However, a standard limitation of these works is their evaluation: they operate entirely in high-precision (FP16/FP32), ignoring the down-stream compression step that is mandatory for practical deployment.

\subsection{Post-Training Quantization (PTQ)}
Post-Training Quantization (PTQ) is a critical step for deploying models on resource-constrained hardware \cite{Gholami2021Survey,Nagel2021DFQ}. Unlike Quantization-Aware Training (QAT), which requires expensive training pipelines and full datasets, PTQ compresses models post-hoc using a tiny calibration set. Prominent methods include GPTQ \cite{Frantar2022GPTQ}, AWQ \cite{Lin2023AWQ}, SmoothQuant \cite{Xiao2023SmoothQuant}, and LLM.int8() \cite{Dettmers2022LLMint8}. Despite their successes, PTQ studies are typically confined to single-task, monolithic models. How PTQ interacts with merged models—where different tasks have been blended together in weight space—remains poorly understood.

\subsection{Methodological Audits of Model Fusion}
The field has seen several attempts to study model merging and quantization. Q-Merge \cite{MergingSub5} proposed a quantization-aware merging scheme using Straight-Through Estimators (STE). However, subsequent work exposing the ``Cross-Schema Generalization Gap'' \cite{MergingSub6} showed that direct optimization of quantization coefficients overfits to the training-time quantization schema and collapses under target schema shifts. Demystifying studies on test-time adaptation \cite{MergingSub7} demonstrated that online adaptations are often fragile under severe shift and that offline validation tuning (OFS-Tune) is far more stable. Furthermore, ZipMerge \cite{MergingSub8} exposed the representation collapse that occurs under joint weight pruning and model merging under extreme domain shifts. 

Despite these progressive investigations, a major methodological blindspot has remained unaddressed: the silent failure of low-rank adapter updates when merged models are re-quantized. While PEFT tutorials claim that adapters can be merged and dequantized freely, we are the first to conduct a multi-axial audit of this ``Re-Quantization Silence.'' We deconstruct the failure mathematically, show how and why standard baselines collapse, and offer robust, scale-aware solutions.